% ===== §21 =====
\section{The Double Interweaving of Drive}
\label{p24:sec:interweaving}

\noindent
The second pillar specifies what, within the vacancy the first established, is actually woven. The general account of Part One held that culture is the exteriorization of drive; the phenomenology of the dyad now shows that in intimate relation the drive so exteriorized has a structure found in no other cultural formation, and that this structure is the mechanism behind the non-transferability the paper has claimed since its engagement with Bourdieu. The structure is a \emph{double interweaving}: the two structures of desire are not exteriorized side by side, in parallel, but are mutually implicated, each entering into the other, so that the culture they generate is the sediment of a single interwoven fantasy rather than the sum of two private ones.

\subsection*{From parallel fantasies to an interwoven one}

Orthodox psychoanalysis writes desire as the subject's relation to its object-cause, the objet petit a around which the drive turns and which veils the constitutive lack. Two people in proximity might, on a thin account, be imagined as two such structures running in parallel: each with a private fantasy, each pursuing a private object-cause, the relation a matter of two desires that happen to have taken each other as their occasion. This is not what an intimate relation is, and the difference is the pillar. In intimate relation each partner comes to occupy, for the other, the place of the object-cause itself, or a part of it: the other is not merely the occasion of a desire whose cause lies elsewhere but is woven into the cause, so that what each desires around is no longer a wholly private object but one in which the other is implicated. The fantasy that structures each one's desire has the other inside it. This is why the desires of two intimate partners cannot be cleanly separated and returned to each as private property: they have grown into each other at the level of the object-cause, and the culture they generate together is the exteriorization not of two adjacent desires but of one interwoven desiring, in which neither partner's drive can be specified without reference to the other.

\subsection*{The private culture as the sediment of the interweaving}

The private forms catalogued earlier are the sediment of exactly this interweaving, and reading them through it explains features that would otherwise be merely charming. A private joke is the crystallization of a shared coordinate of desire: it discharges in laughter because it touches, in a compressed form, the interwoven fantasy the two have built, and it touches it for these two alone because the coordinate it touches exists only in their interweaving. A private name is a condensation-point of the interwoven fantasy: it names the other from within the place the other occupies in one's own object-cause, which is why it would be wrong, even faintly obscene, in a third mouth---the third has no standing in the fantasy the name arises from. A micro-rite is a channel the interwoven desire has cut for itself, a repeated form through which the shared drive circulates. The whole private culture, read this way, is not a collection of shared tastes or agreed conventions but the exteriorized, sedimented form of a fantasy that two people have woven into one, and its every element carries the charge of that interweaving. This is what the outsider cannot reach and the explanation cannot supply: not information about the forms, but a place in the interwoven desire that alone gives the forms their charge.

\subsection*{The mechanism of non-transferability}

With this the paper can state precisely the mechanism of the non-transferability it claimed, against Bourdieu's cultural capital, in Part One. Cultural capital is transferable in principle---accruable, convertible, extractable---because its value is a position in an external field, and a position can, in principle, be occupied by another. The private culture of a couple is non-transferable in principle, and the reason is now exact: its value is lodged in an interweaving of drive that exists only in this relation and cannot be carried out of it. To transfer the culture would require transferring the interwoven fantasy that charges it, and the interwoven fantasy is not a portable content but a structure that came to be only in and as this relation. What can be transferred is the outward form---the sound of the word, the shape of the rite---and what cannot is the charge, because the charge is the interweaving, and the interweaving does not exist outside the two whose desire is interwoven. This is why a couple's private culture, copied by others, is a shell: the others can reproduce the form and cannot reproduce the interwoven desire that made the form a culture rather than a mannerism. Non-transferability is not a contingent privacy that better observation could breach; it is a structural consequence of the fact that the culture's substance is an interweaving of drive proper to one relation.

This can be stated as a thesis, the first of several the paper advances as its own, beyond its revision of Lacan. Call it the \emph{non-transferability thesis}: the degree to which a relational cultural form resists transfer and extraction rises with the depth at which its charge is anchored in the interwoven structure of the participants' drive; and where that anchoring reaches into the unsymbolizable core of the bond, in the register of the Real, what a transfer can carry is the symbolic shell alone, the charge remaining behind because it was never a portable content but the interweaving itself. The thesis is not an identity between transferability and anchoring but a graded relation between them: forms anchored shallowly, in shared convention alone, can be carried off with much of their sense intact, as a borrowed idiom can; forms anchored in the interwoven Real cannot, because the transfer would have to carry an interweaving that exists only in and as the one relation. This is why the non-transferability of an intimate culture admits of degrees, and is most complete for exactly those forms---the private name, the inside joke, the charged silence---that reach deepest into the interwoven drive, and least complete for those that borrow most from the public stock.

\subsection*{The interweaving across the three registers}

The interweaving is not confined to a single register, and marking its span connects this pillar to the fourth. It runs across all three. In the \emph{imaginary}, it is the mutual identification, the way each comes to see themselves in and through the other's image of them, the specular circuit by which two self-images are woven into a shared imago of the couple. In the \emph{symbolic}, it is the private lexicon and the shared references, the coded signs through which the interwoven desire is given a sayable form. In the \emph{real}, it is the unsymbolizable core of the interweaving itself, the ``us'' that no private word quite captures and that is the relational Real of this particular bond. The double interweaving of drive is therefore not merely a symbolic-order phenomenon, a matter of shared signs, but a structure that runs through the imaginary body of identification and down into the real of the unsymbolizable bond. This is why the account of it required the framework of the three registers and could not be given in the terms of the symbolic alone, and it is why the next pillar, having established that the interweaving is what sediments and that it spans the registers, must turn to the distinctive temporality of the sedimentation---the speed at which, in the dense medium of an intimate life, this interwoven drive deposits into the durable forms of a culture.
